\documentclass[aspectratio=169,11pt]{beamer}
\usepackage{../../../shared/theme/esmad_beamer_theme}
\usepackage{amsmath,amssymb}
\usepackage{booktabs}
\usepackage{array}

\title[Instrumental Variables]{Instrumental Variables, Endogeneity, and 2SLS}
\subtitle{Identification from Exogenous Treatment Shifts}
\date{\today}

\newcommand{\E}{\mathbb{E}}
\newcommand{\Cov}{\mathrm{Cov}}
\newcommand{\Var}{\mathrm{Var}}

\begin{document}

\begin{frame}
\titlepage
\end{frame}

\begin{frame}{Learning Goals}
\begin{itemize}
\item explain why endogeneity invalidates the usual OLS causal interpretation;
\item state the relevance, independence, and exclusion conditions for an instrument;
\item derive the Wald estimand and connect it to two-stage least squares;
\item distinguish ATE-style interpretations from Local Average Treatment Effects;
\item diagnose weak instruments and understand their inferential consequences;
\item reason critically about overidentification tests, exclusion restrictions, and invalid instruments.
\end{itemize}
\end{frame}

\begin{frame}{Outline}
\tableofcontents
\end{frame}

\section{Why Instrumental Variables?}

\begin{frame}{The Endogeneity Problem}
Suppose the structural equation is
\[
Y_i=\beta D_i+X_i'\gamma+u_i.
\]

For OLS to identify \(\beta\), we need
\[
\E[u_i\mid D_i,X_i]=0.
\]

Endogeneity means instead
\[
\Cov(D_i,u_i)\neq 0.
\]

\begin{alertblock}{Consequence}
Variation in treatment is partly driven by factors that also affect the outcome, so the regression coefficient mixes causal effect and selection.
\end{alertblock}
\end{frame}

\begin{frame}{Common Sources of Endogeneity}
\begin{itemize}
\item \textbf{Omitted variables}: ability affects education and earnings.
\item \textbf{Simultaneity}: price and quantity are jointly determined.
\item \textbf{Measurement error}: treatment is observed with noise.
\item \textbf{Selection}: people choose treatment based on expected outcomes.
\item \textbf{Policy targeting}: intervention is directed toward units already changing.
\end{itemize}

\begin{block}{IV strategy}
Find a source of treatment variation that is plausibly unrelated to the outcome disturbance except through treatment.
\end{block}
\end{frame}

\section{What Makes an Instrument Valid?}

\begin{frame}{Instrument Relevance}
An instrument \(Z\) must shift treatment:
\[
\Cov(Z,D)\neq 0.
\]

In a first-stage model,
\[
D_i=\pi_0+\pi_1Z_i+X_i'\pi_2+v_i,
\]
we require \(\pi_1\neq 0\).

\begin{alertblock}{Testable, but not sufficient}
A strong first stage only establishes that the instrument predicts treatment. It says nothing by itself about causal validity.
\end{alertblock}
\end{frame}

\begin{frame}{Instrument Independence}
The instrument should be as-good-as-random with respect to the unobserved determinants of the outcome:
\[
Z_i\perp u_i
\]
conditionally on the maintained covariates and design.

\begin{block}{Interpretation}
The assignment mechanism that changes \(Z\) should not select units according to latent outcome determinants.
\end{block}

This is primarily a design and institutional-knowledge claim.
\end{frame}

\begin{frame}{Exclusion Restriction}
The instrument should affect the outcome only through the treatment channel relevant to the estimand.

In structural terms, there should be no direct path
\[
Z\rightarrow Y
\]
after accounting for treatment.

\begin{alertblock}{Usually untestable}
Exclusion is often the most controversial IV assumption and must be defended substantively.
\end{alertblock}
\end{frame}

\begin{frame}{Three Conditions, Three Different Questions}
\begin{tabular}{p{3.1cm}p{4.2cm}p{5cm}}
\toprule
\textbf{Condition} & \textbf{Question} & \textbf{Evidence} \\
\midrule
Relevance & Does \(Z\) move \(D\)? & First-stage data \\
Independence & Is \(Z\) unrelated to latent outcome causes? & Design / institutional knowledge \\
Exclusion & Does \(Z\) affect \(Y\) only through \(D\)? & Scientific argument, falsification \\
\bottomrule
\end{tabular}
\end{frame}

\section{The Wald Estimator}

\begin{frame}{Binary Instrument, Binary Treatment}
For binary \(Z\), a simple IV estimand is
\[
\widehat{\beta}_{Wald}
=
\frac{
\E[Y\mid Z=1]-\E[Y\mid Z=0]
}{
\E[D\mid Z=1]-\E[D\mid Z=0]
}.
\]

The numerator is the reduced-form effect of \(Z\) on \(Y\).

The denominator is the first-stage effect of \(Z\) on \(D\).
\end{frame}

\begin{frame}{Reduced Form and First Stage}
First stage:
\[
D_i=\pi_0+\pi_1Z_i+v_i.
\]

Reduced form:
\[
Y_i=\rho_0+\rho_1Z_i+e_i.
\]

With one instrument and one endogenous treatment,
\[
\widehat{\beta}_{IV}
=
\frac{\widehat{\rho}_1}{\widehat{\pi}_1}.
\]

\begin{block}{Interpretation}
The causal effect is the outcome change induced by the instrument divided by the treatment change induced by the instrument.
\end{block}
\end{frame}

\section{Two-Stage Least Squares}

\begin{frame}{First Stage}
Estimate
\[
D_i=\pi_0+Z_i'\pi_1+X_i'\pi_2+v_i.
\]

This produces predicted treatment
\[
\widehat D_i.
\]

Only the component of treatment variation explained by instruments and included exogenous covariates is retained.
\end{frame}

\begin{frame}{Second Stage}
Estimate
\[
Y_i=\beta_0+\beta_1\widehat D_i+X_i'\beta_2+\varepsilon_i.
\]

\begin{alertblock}{Important implementation point}
Do not literally run two unrelated OLS regressions and use ordinary second-stage standard errors. Proper 2SLS software accounts for first-stage estimation in the variance calculation.
\end{alertblock}
\end{frame}

\begin{frame}{Projection Interpretation}
Let \(P_Z\) denote projection onto the instrument space.

Then
\[
\widehat D=P_ZD.
\]

2SLS estimates the relationship between \(Y\) and the part of \(D\) that lies in the span of exogenous instruments.

\begin{block}{Core idea}
The estimator discards endogenous treatment variation not explained by the instrument.
\end{block}
\end{frame}

\section{Local Average Treatment Effects}

\begin{frame}{Treatment-Response Types}
With binary instrument and treatment, units may be:
\begin{itemize}
\item \textbf{compliers}: take treatment when encouraged by \(Z\);
\item \textbf{always-takers}: treated regardless of \(Z\);
\item \textbf{never-takers}: untreated regardless of \(Z\);
\item \textbf{defiers}: move opposite to the instrument.
\end{itemize}
\end{frame}

\begin{frame}{Monotonicity}
A common assumption rules out defiers:
\[
D_i(1)\geq D_i(0)
\]
for all units.

Under standard IV assumptions and monotonicity, the Wald estimand identifies
\[
LATE
=
\E[Y_i(1)-Y_i(0)\mid \text{compliers}].
\]

\begin{alertblock}{Not necessarily the ATE}
The effect applies to units whose treatment status is changed by the instrument.
\end{alertblock}
\end{frame}

\begin{frame}{Why the LATE Population Matters}
Different instruments may affect different margins of treatment.

Therefore two valid instruments can estimate different causal effects if treatment effects are heterogeneous.

\begin{block}{Example}
Distance to college may affect marginal students differently from tuition subsidies. Both can shift education, but for different groups.
\end{block}
\end{frame}

\section{Weak Instruments}

\begin{frame}{What Is a Weak Instrument?}
A weak instrument has only limited explanatory power for the endogenous treatment.

If the first stage is weak:
\begin{itemize}
\item 2SLS can be severely biased in finite samples;
\item estimates become unstable;
\item conventional normal approximations can be poor;
\item confidence intervals can be misleading.
\end{itemize}
\end{frame}

\begin{frame}{First-Stage Diagnostics}
Common diagnostics include:
\begin{itemize}
\item first-stage coefficient estimates;
\item partial \(R^2\);
\item F-type statistics for excluded instruments;
\item weak-identification robust tests and confidence sets.
\end{itemize}

\begin{alertblock}{Avoid rigid folklore}
Rules such as “F greater than 10” are rough heuristics, not universal guarantees of reliable inference.
\end{alertblock}
\end{frame}

\begin{frame}{Why Weak Instruments Are Dangerous}
With weak relevance,
\[
\widehat{\beta}_{IV}
=
\frac{\widehat{\rho}_1}{\widehat{\pi}_1}
\]
has a noisy denominator.

Small changes in the first stage can then generate very large changes in the IV estimate.

\begin{block}{Practical implication}
Report weak-instrument diagnostics alongside the causal estimate, not as an afterthought.
\end{block}
\end{frame}

\section{Multiple Instruments}

\begin{frame}{Overidentified Models}
If there are more instruments than endogenous regressors, the model is overidentified.

This can improve precision, but every instrument must satisfy the required assumptions.

\begin{alertblock}{More instruments are not automatically better}
Many weak or dubious instruments can worsen finite-sample behavior and complicate interpretation.
\end{alertblock}
\end{frame}

\begin{frame}{Overidentification Tests}
Sargan/Hansen-style tests examine whether different instrument-implied moment conditions are mutually compatible.

A rejection may signal invalid instruments or model misspecification.

A non-rejection does not prove validity.

\begin{block}{Reason}
If multiple instruments violate exclusion in similar ways, they may still appear mutually consistent.
\end{block}
\end{frame}

\section{Instrument Validity in Practice}

\begin{frame}{Threats to the Exclusion Restriction}
Suppose \(Z\) changes treatment \(D\), but also:
\begin{itemize}
\item changes another behavior that affects \(Y\);
\item shifts resources directly;
\item changes expectations;
\item alters peer or spillover environments.
\end{itemize}

Then \(Z\) may have a direct or indirect path to \(Y\) outside the treatment channel.
\end{frame}

\begin{frame}{Falsification Strategies}
Useful checks can include:
\begin{itemize}
\item balance of pre-treatment covariates across instrument values;
\item effects of \(Z\) on outcomes that should not respond;
\item reduced-form timing checks;
\item alternative samples where the treatment channel should be absent;
\item sensitivity analysis to plausible exclusion violations.
\end{itemize}

\begin{alertblock}{Scope}
Falsification can expose some invalid instruments; it cannot certify exclusion.
\end{alertblock}
\end{frame}

\section{Applied Reasoning}

\begin{frame}{Example: Price Endogeneity}
Suppose we want demand elasticity:
\[
Q_i=\beta P_i+u_i.
\]

Price \(P_i\) is endogenous because firms set prices partly in response to demand shocks contained in \(u_i\).

A potential instrument is a cost shifter \(Z_i\) that affects price but is plausibly unrelated to demand except through price.

\begin{block}{Question}
Does the cost shifter also affect availability, quality, or consumer expectations? If yes, exclusion may fail.
\end{block}
\end{frame}

\begin{frame}{Example: Regional Pricing Policy}
Suppose a policy changes prices in some regions because of an administrative rule.

A plausible IV workflow is:
\begin{enumerate}
\item explain the rule generating exogenous variation;
\item quantify the first-stage effect on actual prices;
\item inspect whether the rule changes other regional conditions;
\item estimate the reduced form;
\item estimate 2SLS with appropriate uncertainty;
\item state the LATE population affected by the rule.
\end{enumerate}
\end{frame}

\begin{frame}{A Credible IV Workflow}
\begin{enumerate}
\item State the endogenous treatment and the target estimand.
\item Explain precisely why OLS is biased.
\item Describe the instrument assignment mechanism.
\item Defend relevance, independence, and exclusion separately.
\item Report the first stage and reduced form.
\item Use proper 2SLS inference.
\item Diagnose weak identification.
\item Define the complier population when using LATE logic.
\item Run falsification and sensitivity analyses.
\end{enumerate}
\end{frame}

\begin{frame}{Common Failure Modes}
\begin{itemize}
\item treating a strong correlation with treatment as proof of validity;
\item ignoring direct instrument effects on the outcome;
\item calling a LATE estimate an ATE;
\item using conventional inference with weak instruments;
\item adding many weak instruments for apparent precision;
\item interpreting overidentification non-rejection as validation;
\item failing to explain the substantive assignment mechanism.
\end{itemize}
\end{frame}

\begin{frame}{Synthesis: What Makes IV Credible?}
\begin{itemize}
\item IV replaces endogenous treatment variation with instrument-induced variation.
\item Relevance is necessary but not enough.
\item Independence and exclusion come from design and scientific reasoning.
\item 2SLS is the estimator; the instrument assignment mechanism is the identification strategy.
\item With heterogeneous effects, the estimand may be a LATE for compliers.
\item Weak instruments create bias and unreliable conventional inference.
\item Strong IV analyses make first stage, reduced form, target population, and failure modes explicit.
\end{itemize}
\end{frame}

\begin{frame}{Selected References}
\small
\begin{itemize}
\item Angrist, J. D., Imbens, G. W., and Rubin, D. B. (1996). Identification of Causal Effects Using Instrumental Variables. \textit{Journal of the American Statistical Association}, 91(434), 444--455.
\item Angrist, J. D., and Pischke, J.-S. (2009). \textit{Mostly Harmless Econometrics}. Princeton University Press.
\item Imbens, G. W., and Angrist, J. D. (1994). Identification and Estimation of Local Average Treatment Effects. \textit{Econometrica}, 62(2), 467--475.
\item Staiger, D., and Stock, J. H. (1997). Instrumental Variables Regression with Weak Instruments. \textit{Econometrica}, 65(3), 557--586.
\item Stock, J. H., Wright, J. H., and Yogo, M. (2002). A Survey of Weak Instruments and Weak Identification in Generalized Method of Moments. \textit{Journal of Business and Economic Statistics}, 20(4), 518--529.
\item Wooldridge, J. M. (2010). \textit{Econometric Analysis of Cross Section and Panel Data}, 2nd ed. MIT Press.
\end{itemize}
\end{frame}

\contactslide

\end{document}
